% !TEX TS-program = pdflatex
% !TEX encoding = UTF-8 Unicode
\documentclass[ngerman,11pt]{article}
\usepackage[utf8]{inputenc}
\usepackage[T1]{fontenc}
\usepackage[ngerman]{babel}
\usepackage[a4paper, total={16cm, 25cm}]{geometry}
\usepackage{amsmath}
\usepackage{amsfonts}
\usepackage{amssymb}
\usepackage{multicol}
\usepackage{tcolorbox}
\usepackage{cancel}
\usepackage{enumitem}
\usepackage{hyperref}
\usepackage{array}
\usepackage{listings}
\usepackage{xcolor}

% Listings-Konfiguration
\definecolor{codebg}{RGB}{245,245,245}
\definecolor{filebg}{RGB}{252,248,230}
\definecolor{codekey}{RGB}{0,0,180}
\definecolor{codecom}{RGB}{90,120,90}
\definecolor{codestr}{RGB}{160,60,60}

\lstdefinestyle{bash}{
    language=bash,
    basicstyle=\ttfamily\small,
    backgroundcolor=\color{codebg},
    keywordstyle=\color{codekey}\bfseries,
    commentstyle=\color{codecom}\itshape,
    stringstyle=\color{codestr},
    showstringspaces=false,
    breaklines=true,
    frame=single,
    framesep=4pt,
    xleftmargin=4pt,
    morekeywords={export,echo,env,printenv,source,unset,cat,nano,vim},
    captionpos=t,
    abovecaptionskip=2pt,
    belowcaptionskip=2pt
}

\lstdefinestyle{bashfile}{
    language=bash,
    basicstyle=\ttfamily\small,
    backgroundcolor=\color{filebg},
    keywordstyle=\color{codekey}\bfseries,
    commentstyle=\color{codecom}\itshape,
    stringstyle=\color{codestr},
    showstringspaces=false,
    breaklines=true,
    frame=single,
    framesep=4pt,
    xleftmargin=4pt,
    morekeywords={export},
    captionpos=t,
    abovecaptionskip=2pt,
    belowcaptionskip=2pt
}

\lstdefinestyle{python}{
    language=Python,
    basicstyle=\ttfamily\small,
    backgroundcolor=\color{codebg},
    keywordstyle=\color{codekey}\bfseries,
    commentstyle=\color{codecom}\itshape,
    stringstyle=\color{codestr},
    showstringspaces=false,
    breaklines=true,
    frame=single,
    framesep=4pt,
    xleftmargin=4pt
}

\begin{document}
\hrule
\begin{center}
{\large\bf Umgebungsvariablen unter Linux --- API-Keys sicher verwalten}\\[-0mm]
Studienkolleg der TU Berlin\hfill Till Zoppke\\
Informatik / T-Kurs \hfill \today\\[2mm]
\end{center}
\hrule
\bigskip

% Einleitung
\noindent Aus Python kennen Sie Variablen: Sie speichern Werte unter einem Namen innerhalb eines Programms. Sobald das Programm endet, sind die Variablen weg. Das Betriebssystem selbst kennt aber auch Variablen --- sogenannte \emph{Umgebungsvariablen} (\emph{environment variables}). Diese leben außerhalb Ihres Python-Skripts und können von jedem Programm gelesen werden, das Sie starten. In diesem Arbeitsblatt lernen Sie, Umgebungsvariablen in der \texttt{bash} zu setzen, dauerhaft zu speichern und aus Python heraus zu verwenden --- konkret am Beispiel eines API-Keys für \textbf{Chat-AI} der GWDG Academic Cloud.

\begin{tcolorbox}[title=Definition: Umgebungsvariable]
Eine \textbf{Umgebungsvariable} ist ein Name-Wert-Paar, das die Shell (bei uns \texttt{bash}) und alle von ihr gestarteten Programme lesen können. Konvention: Namen werden \textsc{GROSSGESCHRIEBEN}, Wörter mit Unterstrich getrennt, z.\,B. \texttt{HOME}, \texttt{CHATAI\_API\_KEY}.
\end{tcolorbox}
\bigskip

\section*{1 Variablen in der Shell setzen und anzeigen}

\subsection*{Temporär setzen (nur für die aktuelle Shell-Sitzung)}

\begin{lstlisting}[style=bash,caption={Terminal}]
# Variable setzen: export macht sie auch fuer Unterprozesse sichtbar
export MEIN_NAME="Alex"

# Wert anzeigen: das $-Zeichen ersetzt den Namen durch den Wert
echo $MEIN_NAME

# Alle Umgebungsvariablen auflisten
env
\end{lstlisting}

\subsection*{Bereits vorhandene Systemvariablen}

\begin{center}
\begin{tabular}{|l|p{11cm}|}
\hline
\textbf{Variable} & \textbf{Bedeutung} \\
\hline
\texttt{HOME}  & Pfad zum Home-Verzeichnis, z.\,B. \texttt{/home/tu-berlin.de/alex} \\
\texttt{USER}  & Benutzername des angemeldeten Kontos \\
\texttt{PATH}  & Liste von Verzeichnissen, in denen die Shell nach ausführbaren Programmen sucht \\
\texttt{SHELL} & Pfad zur aktuellen Login-Shell, z.\,B. \texttt{/bin/bash} \\
\texttt{LANG}  & Spracheinstellung, z.\,B. \texttt{de\_DE.UTF-8} \\
\texttt{PWD}   & Aktuelles Arbeitsverzeichnis \\
\hline
\end{tabular}
\end{center}

\section*{2 Variablen persistent machen --- Anwendungsfall Chat-AI}

Beim Schließen des Terminals verschwindet eine mit \texttt{export} gesetzte Variable. Damit sie bei jedem neuen Terminal automatisch wieder geladen wird, trägt man sie in die Konfigurationsdatei \texttt{\textasciitilde/.bashrc} ein. Diese wird bei jedem Start einer interaktiven \texttt{bash} ausgeführt.

\smallskip
Wir arbeiten mit \textbf{Chat-AI} der \emph{GWDG Academic Cloud}. Den persönlichen API-Key haben Sie per E-Mail erhalten. \textbf{Dieser Key gehört nicht in den Quellcode}, er könnte sonst durch Weitergabe der Datei kompromittiert werden. Stattdessen legen Sie ihn als Umgebungsvariable ab.

\smallskip
\noindent \textbf{Schritt 1:} Konfigurationsdatei im Terminal mit \texttt{vim} öffnen.

\begin{lstlisting}[style=bash]
vim ~/.bashrc
\end{lstlisting}

\noindent \textbf{Schritt 2:} In \texttt{vim} mit \texttt{i} in den Einfügemodus wechseln und am Dateiende folgende Zeile ergänzen:

\begin{lstlisting}[style=bashfile]
export CHATAI_API_KEY="sk-hier-steht-ihr-echter-schluessel"
\end{lstlisting}

\noindent \textbf{Schritt 3:} Mit \texttt{Esc} den Einfügemodus verlassen, dann \texttt{:wq} eingeben und \texttt{Enter} drücken --- damit wird gespeichert (\emph{write}) und \texttt{vim} beendet (\emph{quit}).

\smallskip
\noindent \textbf{Schritt 4:} Datei neu einlesen und prüfen, ob der Wert gesetzt ist.

\begin{lstlisting}[style=bash]
source ~/.bashrc
echo $CHATAI_API_KEY
\end{lstlisting}

\section*{3 Key aus Python heraus verwenden}

Das Python-Skript liest den Key aus der Umgebung und schickt einen HTTP-Request an den Chat-AI-Endpoint. Der Key steht \emph{nicht} im Quelltext.

\begin{lstlisting}[style=python]
import json
import os
from urllib.request import Request, urlopen

api_key = os.environ.get("CHATAI_API_KEY")
if api_key is None:
    raise RuntimeError("CHATAI_API_KEY ist nicht gesetzt.")

url = "https://chat-ai.academiccloud.de/v1/chat/completions"
headers = {
    "Authorization": f"Bearer {api_key}",
    "Content-Type": "application/json",
}
payload = {
    "model": "meta-llama-3.1-8b-instruct",
    "messages": [
        {
            "role": "user",
            "content": "Hallo! Bitte stelle dich kurz vor.",
        }
    ],
}

request = Request(
    url,
    data=json.dumps(payload).encode("utf-8"),
    headers=headers,
    method="POST",
)

with urlopen(request) as response:
    data = json.load(response)

print(data["choices"][0]["message"]["content"])
\end{lstlisting}


\end{document}
